\section{秦昭襄王：伊阙之后，如何把胜利变成次序}

秦武王猝卒时，王位落到年轻的异母弟手中。前307年的\eventanchor{WQ-0307-QIN-SUCCESSION}{秦·昭襄王即位}，并不只是年表上一笔继承：宣太后和魏冉所连结的宫廷、外戚与将领网络，先替新王稳住了朝廷，也使军政大权在相当长的时间里不能只用“王命”二字来概括。昭襄王的漫长在位由此开始。他所继承的，是商鞅以来能够征发户口、粮食和兵员的秦国；他所必须学会的，则是如何在太后、相国、宗室和战将各有资源的格局里，使这些资源服从一条较持久的东进路线。

伊水与洛水会合处的峡口，正是这种路线的一次试验场。前293年的\eventref{WQ-0293-YIQUE}{战国·伊阙之战}中，白起所部击溃韩、魏联军，韩魏面向洛阳和函谷关的屏障因而更薄。传世史书把斩获写得极多，又把战果归于白起一人的临阵才能；可以较稳妥把握的是秦军在这一通道取得决定性优势，至于杀伤数字、战斗展开的日次和各军具体位置，仍主要来自后出的叙事，不宜当作可逐项核验的战报。峡口、城邑和渡口比华丽的战功更能说明此战的分量：它让秦更容易把关中仓廪和兵员送到东面，也迫使韩、魏在一条越来越狭窄的防线上安排援军。

但伊阙没有使秦从此只须向东推行。齐、楚仍足以牵制，韩、魏、赵的临时合从也会令函谷关以外的胜利变得昂贵。魏冉经营的封邑、宾客和军功关系，曾为这种扩张提供可调度的人与物；同一条件也留下王权如何约束相权的难题。结构性原因不在于秦廷忽然有了一位能看穿天下的谋士，而在于连续战争已经使每一次出兵都同时成为对粮道、盟约和宫廷指挥权的考验。

范雎进入秦廷，正落在这道考题最紧的时候。关于他的魏国遭遇、改名入秦以及荐引经过，《史记》所写极具人物故事的张力，《战国策》的说辞又有不同层次；这些材料能够说明一个外来游士如何被安置进秦国权力中心，不能让后人把每一段受辱、密谈和报复复原为现场。前266年前后的\eventanchor{WQ-0266-FANJU}{秦·范雎入秦与远交近攻}，较可靠的核心是昭襄王借此调整权力，魏冉终于失势，范雎居相；人事转折所改变的，不只是一个职位，而是宫廷中谁能够把邻国的远近、战争的成本和王的名义连成一项政策。

\begin{textwindow}[《史记》中的一句方略]
“王不如远交而近攻，得寸则王之寸也，得尺亦王之尺也。”

这句见于\sourcebook{史记·范雎蔡泽列传}，后来成为解释秦国外交的钥匙。它以排比和断语组织成一场成功游说，很可能已经受司马迁叙事的剪裁；《战国策》同类说辞亦不宜逐字视为秦廷记录。可是，它准确点出了可由战争地理检验的选择：远方大国即使暂时修好，也未必能替秦守住新得城邑；韩、魏、赵等近邻的土地一旦夺取，才较可能同既有县邑、道路和征发网络相连。“远交近攻”因此不是包治天下的口号，而是把外交的暂时让步同近处的持续吞并配合起来的计算。
\end{textwindow}

这套计算并未使赵自动退让。赵在太行以东的城邑、北方扩张所得的人口以及通向邯郸的道路，使它成为秦在中原北缘最难绕开的对手；秦军的推进也一再遇到抵抗。到前262年，秦攻取野王，韩国上党同本国腹地的交通被截断。上党郡守不愿降秦而求援于赵，赵廷接收上党，于是\eventanchor{WQ-0262-SHANGDANG}{战国·上党归赵}把原先的韩秦战事转成秦赵之间的正面争夺。这里没有一方被动“卷入”：韩国地方官在降与求援之间作出选择，赵要权衡得到山地城邑和承受秦军压力的代价，秦则不能容忍东出通道旁出现一个由赵接手的敌对据点。

直接后果是前线沿太行山地延长，双方都得把粮秣、役夫、守城兵和朝廷的耐心运往山口。中期影响才是\eventref{WQ-0260-CHANGPING}{战国·长平之战}的形成：它不是范雎一句话或白起一场胜利必然推出的终局，而是近攻的战略偏好、上党归属的触发因素、秦赵各自可动员的资源与盟国难以及时协同共同叠加的结果。

\begin{sourcenote}[从人物传记到战争地理]
\sourcebook{资治通鉴}为这一段提供了可连贯对读的编年脊柱；《史记·秦本纪》《白起王翦列传》保存伊阙及秦军推进的基本线索，《范雎蔡泽列传》保存范雎与魏冉更替的中心叙事，《赵世家》则从赵的处境叙述上党。现代研究常以伊阙、野王、上党与太行通道的地理，检验“远交近攻”在粮运和占领上的可行性；同时也提醒我们，范雎游说的完整文本、白起战果的巨数和上党决策的动机，都经过秦汉以后的文本传统塑形。
\end{sourcenote}

\begin{turningpoint}[制度得失：胜利必须接上道路]
昭襄王时期的秦廷把东进的军政资源逐步收束到王权能够主导的次序中，缓解了相权、外戚与军功集团各自牵动外交的风险；代价是近邻国家和边地社会承受了更连续的征发、筑城与战争压力。伊阙打开的不是一条自动通往统一的直线，范雎所概括的也不是无需代价的公式。秦真正获得的，是把一块近处土地、一道山口和下一批军粮接进既有网络的能力；上党危机说明，这种能力一旦同赵的资源和判断相撞，胜利同样可能被拖进一场谁也难以轻易收束的持久战。
\end{turningpoint}
